\chapter{Огляд існуючих методів автоматичної ідентифікації
музичних інструментів}
\label{chap:review}

Автоматична ідентифікація музичних інструментів в аудіозаписах є
складною науково-практичною задачею, що перебуває на перетині теорії
музики, цифрової обробки сигналів і машинного навчання. Тембр
інструментів визначає унікальний характер їх звучання та є ключовою
характеристикою для ідентифікації. На відміну від висоти звуку чи
ритму, тембральні властивості визначаються складним поєднанням
спектральних і часових характеристик аудіосигналу.

Складність автоматичної ідентифікації зумовлена кількома чинниками.
По-перше, у реальних записах кілька інструментів звучать одночасно,
утворюючи акустичну суміш, з якої необхідно виокремити окремі тембри.
По-друге, один і той самий інструмент може звучати по-різному залежно
від виконавця, прийому гри та акустики приміщення. По-третє, задача є
multi-label за своєю природою: кожному запису може відповідати довільна
підмножина інструментальних міток, що суттєво ускладнює вибір і
навчання моделі.

У науковій літературі запропоновано різні підходи до автоматичної
ідентифікації музичних інструментів. У цьому розділі розглянуто наявні
підходи як існуючі рішення у галузі музичної інформатики з метою
обґрунтування вибору методів, що застосовуються у даній роботі.

% ------------------------------------------------------------------
\section{Тембральні характеристики музичних інструментів}
\label{sec:timbre_theory}

Тембр є комплексною перцептивною характеристикою звуку, що відрізняє
один інструмент від іншого за однакової висоти та гучності. Фізично
тембр визначається спектральною огинаючою сигналу, часовою динамікою
амплітуди (атака, затримка, утримання, спад --- ADSR) та наявністю
шумових складових~\cite{muller2015}.

Спектральні характеристики інструментів суттєво відрізняються:
\begin{itemize}
    \item струнні інструменти (скрипка, гітара) мають багатий
    обертоновий спектр з характерними резонансними частотами;
    \item духові інструменти характеризуються специфічними формантними
    областями;
    \item ударні інструменти мають шумоподібний спектр з короткою
    атакою та швидким затуханням.
\end{itemize}

Мел-спектрограма є стандартним представленням для аналізу тембру,
оскільки логарифмічна мел-шкала частот відповідає слуховому
сприйняттю людини та добре відображає спектральні відмінності між
інструментами~\cite{piston1987}.

% ------------------------------------------------------------------
\section{Rule-based підходи до ідентифікації інструментів}
\label{sec:rule_based}

Ранні підходи до ідентифікації інструментів ґрунтувалися на
формалізації знань про тембральні характеристики кожного інструменту.
Вони передбачали ручне конструювання ознак --- центроїд спектру,
спектральний розкид, нульові переходи --- та їх порівняння з
наперед визначеними шаблонами або порогами.

Rule-based методи добре працюють для ізольованих інструментів у
контрольованих умовах. Проте вони є чутливими до акустичних умов
запису, не враховують варіацій виконавської техніки та погано
узагальнюються на суміші інструментів. У задачі multi-label
класифікації такі підходи потребують окремого ансамблю детекторів
для кожного інструменту, що ускладнює реалізацію.

% ------------------------------------------------------------------
\section{Класичні методи машинного навчання}
\label{sec:classical_ml}

Застосування класичних методів машинного навчання --- Gaussian Mixture
Models (GMM), метод опорних векторів (SVM), k-найближчих сусідів (kNN)
--- суттєво підвищило якість ідентифікації порівняно з rule-based
підходами. Такі методи навчаються автоматично розрізняти інструменти на
основі наборів спектральних і часових ознак~\cite{ellis2007}.

Для multi-label задачі класичні ML-методи адаптуються через підходи
«один проти одного» або «один проти всіх», де для кожного інструменту
навчається окремий бінарний класифікатор. Серед ознак
широко застосовуються MFCC, спектральний центроїд, спектральний потік
та мел-спектральні коефіцієнти.

Обмеженням класичних ML-методів є те, що вони обробляють кожен фрагмент
незалежно та не здатні враховувати часовий контекст звучання
інструментів. Крім того, якість результатів суттєво залежить від
ручного конструювання ознак.

% ------------------------------------------------------------------
\section{Методи цифрової обробки сигналів}
\label{sec:dsp}

Ключовим етапом будь-якого підходу до ідентифікації інструментів є
виділення інформативних ознак з аудіосигналу. Методи цифрової обробки
забезпечують перехід від сирого аудіосигналу до компактного
представлення, що зберігає тембральну інформацію.

Перетворення Фур'є короткого вікна (STFT) є базовим інструментом
часово-частотного аналізу. На його основі будується мел-спектрограма
--- часово-частотне представлення з мел-шкалою частот, що
відповідає особливостям людського слуху. Мел-частотні кепстральні
коефіцієнти (MFCC) є компактним описом спектральної огинаючої та
широко використовуються для представлення тембру~\cite{harte2006}.

Порівняно з хроматичними ознаками, що застосовуються у гармонічному
аналізі, мел-спектрограми та MFCC є більш інформативними для задачі
ідентифікації тембру, оскільки зберігають інформацію про весь
спектр сигналу, а не лише про клас висоти звуку.

% ------------------------------------------------------------------
\section{Згорткові нейронні мережі (CNN)}
\label{sec:cnn}

Згорткові нейронні мережі продемонстрували високу ефективність у
задачах аналізу зображень і були успішно адаптовані для обробки
спектральних представлень аудіосигналу. Мел-спектрограма природно
інтерпретується як двовимірний образ, де осі відповідають часу та
частоті, що дозволяє застосовувати стандартні CNN-архітектuri для
виявлення тембральних патернів~\cite{korzeniowski2017}.

Двовимірні згорткові фільтри здатні автоматично виявляти локальні
часово-частотні закономірності, характерні для конкретних
інструментів: атакові транзієнти ударних, гармонійні смуги струнних,
формантні структури духових. Операція пулінгу забезпечує інваріантність
до малих зсувів у часі та частоті, підвищуючи стійкість до виконавських
варіацій.

% ------------------------------------------------------------------
\section{Рекурентні нейронні мережі та гібридні архітектури}
\label{sec:rnn}

Рекурентні нейронні мережі, зокрема BiLSTM, дозволяють враховувати
часовий контекст звучання інструментів~\cite{hochreiter1997}. Для задачі
ідентифікації інструментів це важливо, оскільки тембральний образ
інструменту проявляється протягом певного часового проміжку: атака,
утримання та спад є складовими, що разом визначають впізнаваний тембр.

Гібридні архітектури CNN~+~LSTM (або CNN~+~BiLSTM) поєднують
можливості CNN для виявлення локальних тембральних патернів зі
здатністю рекурентних мереж моделювати часові залежності. Такий підхід
є природним для задачі ідентифікації інструментів, де необхідно
враховувати як спектральну структуру (тембр), так і динаміку звучання
у часі~\cite{bock2012}.

% ------------------------------------------------------------------
\section{Огляд існуючих наборів даних та інструментів оцінювання}
\label{sec:datasets}

Для об'єктивного порівняння методів автоматичної ідентифікації
інструментів застосовуються стандартизовані набори даних з розміченими
інструментальними анотаціями.

\subsection{IRMAS}
\label{subsec:irmas}

Датасет IRMAS (Instrument Recognition in Musical Audio Signals)~\cite{irmas2012}
містить близько 6000 тризекундних фрагментів аудіозаписів з анотаціями
для 11 класів інструментів: гітара, електрогітара, орган, фортепіано,
труба, скрипка, гітара-флейта, флейта, ударні, контрабас та голос.
Набір розроблено спеціально для задачі ідентифікації інструментів у
полімодальних записах і є стандартним бенчмарком у галузі.

\subsection{OpenMIC-2018}
\label{subsec:openmic}

OpenMIC-2018~\cite{openmic2018} містить близько 20\,000 десятисекундних
кліпів із слабкими анотаціями для 20 класів інструментів. Слабка
анотація означає, що мітки отримано частково за допомогою краудсорсингу
та не гарантують повноти: відсутність мітки не означає відсутності
інструменту у записі. Це робить набір більш реалістичним, проте
потребує спеціальних підходів до навчання.

\subsection{Метрики оцінювання}
\label{subsec:metrics}

Для оцінювання якості методів multi-label класифікації застосовуються
метрики, що враховують помилки на рівні окремих міток. F1-macro
обчислюється як незважене середнє F1-міри по всіх класах, F1-micro
--- як F1-міра по об'єднаній множині всіх міток. Середня точність
mAP (mean Average Precision) є стандартною метрикою для задач з
нерівномірним розподілом класів і вимірює якість ранжування передбачених
міток.

% ------------------------------------------------------------------
\section{Висновки до розділу}
\label{sec:ch2_conclusions}

У розділі проведено огляд існуючих методів автоматичної ідентифікації
музичних інструментів в аудіозаписах. Усі розглянуті підходи є наявними
у науковій літературі рішеннями і не є розробками даної роботи.
Узагальнене порівняння розглянутих класів методів наведено у
таблиці~\ref{tab:methods_comparison}.

\begin{table}[H]
\centering
\caption{Порівняння підходів до автоматичної ідентифікації інструментів}
\label{tab:methods_comparison}
\begin{tabular}{|p{3.0cm}|p{4.2cm}|p{4.2cm}|p{2.2cm}|}
\hline
\textbf{Підхід} & \textbf{Переваги} & \textbf{Недоліки} &
\textbf{Точність} \\
\hline
Rule-based &
Інтерпретованість, не потребує навчання &
Чутливість до умов запису, не узагальнюється &
Низька \\
\hline
Класичні ML &
Простота реалізації, менші вимоги до даних &
Не враховує часовий контекст, ручне конструювання ознак &
Помірна \\
\hline
CNN &
Автоматичне виявлення тембральних патернів &
Не враховує довгострокові часові залежності &
Висока \\
\hline
CNN + BiLSTM &
Поєднує тембральний аналіз і часовий контекст &
Вища складність навчання &
Висока \\
\hline
\end{tabular}
\end{table}

віsds Rule-based методи та класичні ML-підходи не здатні ефективно моделювати
тембральну складність реальних полімодальних записів. Виключно CNN-архітектури
не враховують часовий контекст звучання інструментів.

З урахуванням проведеного аналізу для вирішення поставленої задачі
обрано гібридну архітектуру CNN~+~BiLSTM у поєднанні з мел-спектрограмами.
Така комбінація дозволяє виявляти тембральні патерни за допомогою
згорткових шарів та враховувати часовий контекст звучання за допомогою
BiLSTM. Для навчання та оцінювання обрано датасет IRMAS як стандартизований
бенчмарк з детальними інструментальними анотаціями.
